\noindent (b) 试作出 $\mu_{1}$ 和 $\mu_{2}$ 的 $95\%$ 联合 $T^{2}$ 区间及庞弗罗尼区间, 并将这两种区间进行比较. $T^{2}$ 区间是否比庞弗罗尼区间更优越 (如果有优越性的话)?

\noindent (c) 二元正态分布是可行的总体模型吗? 试用 Q-Q 图和散布图来解释.

\medskip
\noindent\mvbold{5.21}\quad 利用\hyperref[tab:1.8]{表 1.8} 中的数据, 构造其各均值的 $95\%$ 庞弗罗尼区间, 以及 $95\%$ 联合 $T^{2}$ 区间, 并对这两个区间进行比较.

\medskip
\noindent\mvbold{5.22}\quad 复制第 6 章\hyperref[tab:6.10]{表 6.10} 中的部分数据 (\hyperref[tab:5.13]{表 5.13}). 这些数据代表用卡车把牛奶从农场运到奶制品厂的各种运输成本. \hyperref[tab:6.10]{表 6.10} 只给出了卡车的前 25 个多元观测值. 第 9 个和第 21 个观测值是由 36 个观测值组成的整个数据集中的两个离群值 (见文献 [2]).

\begin{mvtable}
\phantomsection\label{tab:5.13}\mvtabcap{表 5.13\quad 牛奶运输价格数据}
\centering
{\small
\begin{tabular}{|c|c|c||c|c|c|}
\hline
燃料($x_{1}$) & 维修($x_{2}$) & 资金($x_{3}$) & 燃料($x_{1}$) & 维修($x_{2}$) & 资金($x_{3}$)\\
\hline
$16.44$ & $12.43$ & $11.23$ & $11.11$ & $6.15$ & $7.61$\\
$7.19$ & $2.70$ & $3.92$ & $12.17$ & $14.26$ & $14.39$\\
$9.92$ & $1.35$ & $9.75$ & $10.24$ & $2.59$ & $6.09$\\
$4.24$ & $5.78$ & $7.78$ & $10.18$ & $6.05$ & $12.14$\\
$11.20$ & $5.05$ & $10.67$ & $8.88$ & $2.70$ & $12.23$\\
$14.25$ & $5.78$ & $9.88$ & $12.34$ & $7.73$ & $11.68$\\
$13.50$ & $10.98$ & $10.60$ & $8.51$ & $14.02$ & $12.01$\\
$13.32$ & $14.27$ & $9.45$ & $26.16$ & $17.44$ & $16.89$\\
$29.11$ & $15.09$ & $3.28$ & $12.95$ & $8.24$ & $7.18$\\
$12.68$ & $7.61$ & $10.23$ & $16.93$ & $13.37$ & $17.59$\\
$7.51$ & $5.80$ & $8.13$ & $14.70$ & $10.78$ & $14.58$\\
$9.90$ & $3.63$ & $9.13$ & $10.32$ & $5.16$ & $17.00$\\
$10.25$ & $5.07$ & $10.17$ & & &\\
\hline
\end{tabular}}
\end{mvtable}

\noindent (a) 试作燃料成本、维修成本与资金成本的边缘分布的 Q-Q 图, 并作出各对不同变量观测值之间的三个可能的散布图. 有无明显的离群值? 摒弃离群值后, 重新绘制 Q-Q 图及散布图, 试讨论现在的数据是否服从正态分布.

\noindent (b) 求成本均值的 $95\%$ 庞弗罗尼区间及 $T^{2}$ 区间, 并对它们进行比较.

\medskip
\noindent\mvbold{5.23}\quad 考虑\hyperref[tab:6.13]{表 6.13} 中的第一时段的 30 个古埃及男性头骨观测值.

\noindent (a) 构造头骨的最大宽度, 头骨的高度, 头骨底穴至齿槽的长度和头骨鼻梁高度这些变量的边缘分布的 Q-Q 图. 另外, 构造一个多元观测的 $\chi^{2}$ 图. 这些数据服从正态分布吗? 请解释原因.

\noindent (b) 构造个体头骨尺寸变量的 $95\%$ 庞弗罗尼区间. 另外, 找出它的 $95\%$ $T^{2}$ 区间. 比较这两个区间.

\medskip
\noindent\mvbold{5.24}\quad 利用\hyperref[tab:5.8]{表 5.8} 中的麦迪逊警署数据, 构造 $x_{1}$ 和 $x_{2}$ 的个别 $\bar{X}$ 控制图. 试问这两个过程是否都处在受控状态 (即是否稳定)?

\medskip
\noindent\mvbold{5.25}\quad 参看练习 5.24. 利用 $x_{2}$ 与 $x_{1}$ 的数据, 作出它们的质量控制椭圆及 $T^{2}$ 控制图. 这些二元观测值所代表的过程是否处在受控状态? 试作出评论. 从多元控制图中, 你是否学到了某些在个别 $\bar{X}$ 控制图中没有的东西?

\medskip
\noindent\mvbold{5.26}\quad 利用\hyperref[tab:5.8]{表 5.8} 中的 $x_{1},x_{2}$ 及 $x_{3}$ 的数据构造一个 $T^{2}$ 控制图, 将此图与例 5.10 中的
